\subsection{Thiết kế cơ sở dữ liệu}

Hệ thống sử dụng MongoDB làm cơ sở dữ liệu chính và Supabase Storage/S3-compatible storage làm kho lưu trữ ảnh evidence. MongoDB lưu metadata của người dùng, phiên đăng nhập, phiên xử lý, kết quả nhận dạng và sự kiện được đánh dấu. Supabase Storage không được dùng như một hệ quản trị cơ sở dữ liệu quan hệ, mà chỉ lưu các file ảnh như crop phương tiện, crop biển số và các frame trong track buffer; MongoDB lưu đường dẫn hoặc URL công khai đến các ảnh này.

Cách thiết kế này phù hợp với đặc thù dữ liệu của bài toán ALPR. Mỗi phiên xử lý video sinh ra nhiều phương tiện được theo dõi, mỗi phương tiện lại có nhiều crop biển số theo thời gian. Nếu lưu toàn bộ ảnh nhị phân vào một document, kích thước document dễ tăng nhanh và gây khó khăn cho truy vấn lịch sử. Vì vậy, hệ thống chỉ lưu metadata trong MongoDB, còn ảnh evidence được tách sang object storage. Đối với các dữ liệu nhỏ và thường được truy xuất cùng nhau, ví dụ các frame trong bộ đệm của một record hoặc danh sách phương tiện trong một event, hệ thống sử dụng embedded document để giảm số lần truy vấn.

\subsubsection{Biểu đồ thực thể liên kết logic}

Mặc dù hệ thống sử dụng MongoDB, biểu đồ E-R vẫn được dùng để mô tả quan hệ logic giữa các thực thể chính. Sau đó, các quan hệ này được ánh xạ sang mô hình document bằng hai cách: tham chiếu giữa các collection hoặc nhúng document con trong document cha.

\begin{figure}[H]
    \centering
    \includegraphics[width=0.95\linewidth]{Hinhve/database_er.png}
    \caption{Biểu đồ E-R logic và ánh xạ NoSQL của hệ thống}
    \label{fig:database-er}
\end{figure}

Các thực thể chính trong hệ thống gồm:

\begin{itemize}
    \item \textbf{User}: biểu diễn tài khoản người dùng của hệ thống. Một người dùng có thể tạo nhiều phiên đăng nhập, nhiều phiên xử lý video và nhiều record nhận dạng.
    \item \textbf{AuthSession}: biểu diễn một phiên đăng nhập của người dùng. Thực thể này phục vụ cơ chế xác thực bằng session cookie phía server.
    \item \textbf{RecognitionSession}: biểu diễn một phiên xử lý video hoặc nguồn dữ liệu nhận dạng. Một phiên xử lý có nhiều record biển số tương ứng với các phương tiện được phát hiện và theo dõi.
    \item \textbf{RecognitionRecord}: biểu diễn kết quả nhận dạng cuối cùng của một phương tiện trong một phiên xử lý. Record lưu biển số, độ tin cậy OCR, loại phương tiện, thông tin track và các evidence liên quan.
    \item \textbf{PlateFrame}: biểu diễn một frame crop biển số trong quá trình theo dõi. Đây là document con được nhúng vào \texttt{RecognitionRecord}, không phải collection độc lập.
    \item \textbf{Event}: biểu diễn một sự kiện hoặc cửa sổ thời gian ngắn được người dùng đánh dấu trong chế độ monitor.
    \item \textbf{EventVehicle}: biểu diễn phương tiện được phát hiện trong một event. Đây là document con được nhúng trong \texttt{Event}.
    \item \textbf{EvidenceObject}: biểu diễn các file ảnh được lưu trong Supabase Storage/S3. Thực thể này không phải collection trong MongoDB, mà là tài nguyên ngoài được tham chiếu bằng URL.
\end{itemize}

Quan hệ giữa các thực thể được thiết kế như sau. Một \texttt{User} có nhiều \texttt{AuthSession}; quan hệ này được triển khai bằng trường \texttt{user\_id} trong collection \texttt{auth\_sessions}. Một \texttt{User} cũng có nhiều \texttt{RecognitionSession}; mỗi session lưu \texttt{user\_id} để phục vụ truy vấn lịch sử theo tài khoản. Một \texttt{RecognitionSession} có nhiều \texttt{RecognitionRecord}; mỗi record lưu \texttt{session\_id} và \texttt{user\_id}. Cách lưu bằng reference giúp hệ thống dễ liệt kê phiên gần nhất, sau đó lấy các record thuộc phiên được chọn.

Đối với \texttt{PlateFrame}, hệ thống không tách thành collection riêng. Mỗi \texttt{RecognitionRecord} chứa một \texttt{best\_plate\_frame} và một mảng \texttt{track\_buffer}. Đây là quan hệ ``contains'', trong đó frame chỉ có ý nghĩa trong ngữ cảnh của record chứa nó. Khi người dùng mở chi tiết evidence hoặc xem track buffer, ứng dụng cần lấy cả record và danh sách frame liên quan, do đó nhúng \texttt{PlateFrame} giúp giảm số lần truy vấn.

Tương tự, \texttt{EventVehicle} được nhúng trong \texttt{Event}. Một event thường được xem như một đơn vị phân tích hoàn chỉnh: khoảng thời gian được đánh dấu, trạng thái xử lý, số phương tiện và evidence của từng phương tiện. Vì vậy, nhúng danh sách phương tiện vào event giúp việc đọc chi tiết event đơn giản hơn.

\subsubsection{Ánh xạ sang mô hình MongoDB}

Từ mô hình E-R logic, hệ thống triển khai thành các collection MongoDB như Bảng \ref{tab:mongodb-collections}. Các quan hệ có kích thước tăng trưởng lớn hoặc cần truy vấn độc lập được triển khai bằng reference. Các quan hệ nhỏ, phụ thuộc ngữ cảnh và thường được đọc cùng document cha được triển khai bằng embedded document.

\begin{table}[H]
\centering
\caption{Ánh xạ thực thể sang collection/document MongoDB}
\label{tab:mongodb-collections}
\begin{tabular}{|p{3.4cm}|p{3.6cm}|p{7.2cm}|}
\hline
\textbf{Thực thể logic} & \textbf{Collection/Document} & \textbf{Mô tả thiết kế} \\
\hline
User & \texttt{users} & Collection độc lập. Lưu thông tin tài khoản, email, tên hiển thị, mật khẩu đã băm, vai trò và trạng thái hoạt động. \\
\hline
AuthSession & \texttt{auth\_sessions} & Collection độc lập. Tham chiếu đến \texttt{users} bằng \texttt{user\_id}; dùng để kiểm tra phiên đăng nhập và thu hồi session. \\
\hline
RecognitionSession & \texttt{recognition\_sessions} & Collection độc lập. Lưu metadata của phiên xử lý như file nguồn, trạng thái, tổng số record, số frame đã xử lý, chế độ tiền xử lý và OCR backend. \\
\hline
RecognitionRecord & \texttt{recognition\_records} & Collection độc lập. Tham chiếu đến \texttt{recognition\_sessions} bằng \texttt{session\_id}; lưu kết quả biển số cuối cùng của từng track. \\
\hline
PlateFrame & Embedded document & Được nhúng trong \texttt{RecognitionRecord} dưới dạng \texttt{best\_plate\_frame} và mảng \texttt{track\_buffer}. \\
\hline
Event & \texttt{events} & Collection độc lập. Lưu sự kiện được đánh dấu trong chế độ monitor, gồm nguồn, khoảng thời gian, trạng thái và thời gian xử lý. \\
\hline
EventVehicle & Embedded document & Được nhúng trong \texttt{Event} dưới dạng mảng \texttt{vehicles}. \\
\hline
EvidenceObject & Supabase/S3 object & Không phải collection MongoDB. File ảnh được lưu trong bucket object storage; MongoDB chỉ lưu URL trong các trường \texttt{image\_url}, \texttt{vehicle\_thumbnail\_url}, \texttt{plate\_image\_url}, \texttt{vehicle\_image\_url}. \\
\hline
\end{tabular}
\end{table}

\subsubsection{Thiết kế các collection chính}

\textbf{Collection \texttt{users}.} Collection này lưu thông tin người dùng của hệ thống. Các trường chính gồm \texttt{\_id}, \texttt{email}, \texttt{name}, \texttt{password\_hash}, \texttt{role}, \texttt{is\_active}, \texttt{created\_at} và \texttt{updated\_at}. Trường \texttt{email} được đặt unique index để không có hai tài khoản dùng cùng một email. Hệ thống không lưu mật khẩu gốc, chỉ lưu mật khẩu đã băm.

\textbf{Collection \texttt{auth\_sessions}.} Collection này lưu phiên đăng nhập phía server. Mỗi document gồm \texttt{session\_id}, \texttt{user\_id}, \texttt{expires\_at}, \texttt{revoked}, \texttt{created\_at} và \texttt{updated\_at}. Khi người dùng đăng nhập, backend tạo một \texttt{AuthSession}, sinh JWT chứa \texttt{user\_id} và \texttt{session\_id}, sau đó gửi về trình duyệt dưới dạng HttpOnly cookie. Khi cần xác thực request, backend kiểm tra cookie, giải mã JWT và đối chiếu với document trong \texttt{auth\_sessions}.

\textbf{Collection \texttt{recognition\_sessions}.} Collection này biểu diễn một phiên xử lý video. Các trường quan trọng gồm \texttt{session\_id}, \texttt{user\_id}, \texttt{source\_filename}, \texttt{source\_type}, \texttt{status}, \texttt{total\_records}, \texttt{processed\_frames}, \texttt{preprocess\_mode}, \texttt{ocr\_backend}, \texttt{error\_message}, \texttt{created\_at} và \texttt{updated\_at}. Khi người dùng upload video, backend tạo document session ở trạng thái \texttt{processing}; khi pipeline hoàn tất, session được cập nhật thành \texttt{completed} hoặc \texttt{failed}.

\textbf{Collection \texttt{recognition\_records}.} Collection này lưu kết quả cuối cùng cho từng phương tiện được theo dõi trong một phiên. Các trường định danh gồm \texttt{session\_id}, \texttt{user\_id}, \texttt{track\_id}, \texttt{vehicle\_track\_id} và \texttt{plate\_track\_id}. Các trường kết quả gồm \texttt{vehicle\_class}, \texttt{plate\_text}, \texttt{plate\_text\_confidence}, \texttt{ocr\_vote\_summary}, \texttt{ocr\_method}, \texttt{first\_seen\_frame} và \texttt{last\_seen\_frame}. Các trường evidence gồm \texttt{vehicle\_thumbnail\_url}, \texttt{best\_plate\_frame} và \texttt{track\_buffer}.

Trong \texttt{recognition\_records}, \texttt{best\_plate\_frame} là một document con chứa frame biển số tốt nhất, còn \texttt{track\_buffer} là mảng các \texttt{PlateFrame} được chọn từ quá trình tracking. Mỗi \texttt{PlateFrame} lưu \texttt{frame\_index}, \texttt{quality\_score}, \texttt{image\_url}, \texttt{ocr\_text} và \texttt{ocr\_confidence}. Ảnh crop thật được upload lên object storage, còn \texttt{image\_url} trỏ đến ảnh đó.

\textbf{Collection \texttt{events}.} Collection này lưu kết quả phân tích sự kiện được đánh dấu trong chế độ monitor. Các trường chính gồm \texttt{event\_id}, \texttt{session\_id}, \texttt{source\_type}, \texttt{source\_ref}, \texttt{marked\_at}, \texttt{window\_start\_sec}, \texttt{window\_end\_sec}, \texttt{duration\_sec}, \texttt{status}, \texttt{vehicles}, \texttt{total\_vehicles}, \texttt{processing\_ms}, \texttt{error\_message}, \texttt{created\_at} và \texttt{updated\_at}. Trường \texttt{vehicles} là mảng \texttt{EventVehicle}. Mỗi phần tử lưu biển số, độ tin cậy, loại phương tiện, số frame OCR và URL ảnh evidence.

\subsubsection{Thiết kế lưu trữ evidence bằng Supabase Storage/S3}

Các ảnh evidence được lưu vào bucket \texttt{evidence}. Trong luồng xử lý video thông thường, hệ thống lưu ảnh theo mẫu đường dẫn:

\begin{itemize}
    \item \texttt{\{session\_id\}/track\_\{track\_id\}\_frame\_\{index\}.jpg}: các frame trong track buffer;
    \item \texttt{\{session\_id\}/plate\_\{track\_id\}.jpg}: ảnh biển số tốt nhất;
    \item \texttt{\{session\_id\}/vehicle\_\{track\_id\}.jpg}: ảnh crop phương tiện.
\end{itemize}

Trong luồng monitor event, ảnh evidence được lưu theo mẫu:

\begin{itemize}
    \item \texttt{events/\{event\_id\}/plate\_\{vehicle\_id\}.jpg};
    \item \texttt{events/\{event\_id\}/vehicle\_\{vehicle\_id\}.jpg}.
\end{itemize}

Sau khi upload ảnh, backend lấy public URL và lưu URL đó vào MongoDB. Cách lưu này tách dữ liệu nhị phân khỏi metadata. Giao diện lịch sử và evidence chỉ cần đọc document MongoDB để lấy URL, sau đó hiển thị ảnh trực tiếp từ object storage. Nếu triển khai trong môi trường cần bảo mật evidence chặt hơn, bucket có thể chuyển sang private và backend sẽ sinh signed URL có thời hạn thay vì lưu public URL.

\subsubsection{Thiết kế chỉ mục}

Các chỉ mục được tạo để phục vụ đúng các truy vấn chính của ứng dụng: đăng nhập bằng email, kiểm tra session, liệt kê lịch sử theo người dùng, lấy record theo phiên và lấy event gần nhất. Bảng \ref{tab:mongodb-indexes} liệt kê các chỉ mục chính.

\begin{table}[H]
\centering
\caption{Các chỉ mục MongoDB chính}
\label{tab:mongodb-indexes}
\begin{tabular}{|p{3.8cm}|p{5.2cm}|p{5.2cm}|}
\hline
\textbf{Collection} & \textbf{Chỉ mục} & \textbf{Mục đích} \\
\hline
\texttt{users} & \texttt{email} unique & Đảm bảo email không trùng và tăng tốc đăng nhập. \\
\hline
\texttt{auth\_sessions} & \texttt{session\_id} unique, \texttt{user\_id}, \texttt{expires\_at} & Kiểm tra phiên đăng nhập, truy vết phiên theo người dùng và hỗ trợ xử lý phiên hết hạn. \\
\hline
\texttt{recognition\_sessions} & \texttt{session\_id} unique, \texttt{(user\_id, created\_at)}, \texttt{status}, \texttt{created\_at} & Lấy phiên theo mã, liệt kê lịch sử gần nhất của người dùng và lọc theo trạng thái xử lý. \\
\hline
\texttt{recognition\_records} & \texttt{session\_id}, \texttt{(user\_id, session\_id)}, \texttt{plate\_text}, \texttt{vehicle\_track\_id}, \texttt{plate\_track\_id}, \texttt{vehicle\_class}, \texttt{created\_at}, \texttt{(session\_id, track\_id)} unique & Lấy record của một phiên, lấy record theo track, tìm kiếm biển số, lọc theo mã track phương tiện/biển số và tránh lưu trùng cùng một track trong một session. \\
\hline
\texttt{events} & \texttt{event\_id} unique, \texttt{session\_id}, \texttt{marked\_at}, \texttt{status}, \texttt{source\_type} & Lấy chi tiết event, liệt kê event theo phiên, theo nguồn và theo thời gian đánh dấu. \\
\hline
\end{tabular}
\end{table}

\subsubsection{Luồng truy vấn chính}

Thiết kế collection và chỉ mục được chọn dựa trên các luồng truy vấn chính của ứng dụng. Khi người dùng đăng nhập, hệ thống tìm tài khoản theo \texttt{email}, kiểm tra mật khẩu và tạo document trong \texttt{auth\_sessions}. Khi người dùng mở lịch sử, frontend gọi API lấy danh sách \texttt{recognition\_sessions} theo \texttt{user\_id}, sắp xếp giảm dần theo \texttt{created\_at}. Khi người dùng chọn một phiên, backend lấy các \texttt{recognition\_records} theo \texttt{session\_id} và \texttt{user\_id}. Khi người dùng mở track buffer, backend lấy một record theo cặp \texttt{(session\_id, track\_id)}.

Đối với monitor event, khi người dùng đánh dấu một khoảng thời gian, backend chạy pipeline ALPR trên cửa sổ đó và lưu kết quả vào \texttt{events}. Khi cần xem lại event, ứng dụng truy vấn theo \texttt{event\_id}, \texttt{session\_id} hoặc \texttt{source\_type}. Vì danh sách phương tiện trong event được nhúng trực tiếp vào document event, giao diện có thể hiển thị evidence của sự kiện bằng một lần đọc document chính.

Tóm lại, thiết kế cơ sở dữ liệu của hệ thống kết hợp giữa reference và embedded document. Reference được dùng cho các thực thể lớn, có vòng đời riêng và cần truy vấn độc lập như người dùng, phiên xử lý, record và event. Embedded document được dùng cho dữ liệu phụ thuộc ngữ cảnh như \texttt{PlateFrame} và \texttt{EventVehicle}. Ảnh evidence được tách khỏi MongoDB và lưu trong object storage để giảm kích thước document, trong khi MongoDB vẫn giữ đủ metadata và URL để phục vụ truy vấn, lịch sử và hiển thị kết quả.
